\section{Vectors detected at a prescribed weight}\label{pair:section}

The next step makes the path coding usable inside arbitrary subspaces.  We
must be able to prescribe the next weight of a path and still find a bounded
vector on which its next piece has a substantial value.  All constants in
this section are uniform in the family and in the prescribed weight.

Fix one family.  In a pure family, an admissible block sequence consists of
nonzero finite vectors on its color with successive ordinal supports.  In
the mixed family, its supports are successive in both ordinal and color
order.  A finite block of such a sequence is a nonzero finite linear combination
of its members; successive blocks use successive finite sets of indices.  We use the parameters
$s,r=s',q,L_j,m_j,\theta_j,c_0,B_j,d_j$ of this family from
Lemma~\ref{ch:parameters}.  In particular $1<q<r$, and
$\theta_j=L_j^{1/r}/m_j$.

\begin{proposition}\label{pair:exact}
Let $(z_n)$ be an infinite admissible block sequence in $E$.  Given a tail
of this sequence and an integer $J\ge1$, there are a finite block $y$ in
that tail and a Type~I functional $\phi$ of weight $J$ in the same family
such that
\[
 \|y\|\le500,\qquad \phi(y)\ge\tfrac12.
\]
The support of $\phi$ lies in the ordinal range of $y$, and also in its
color range in the mixed case.
\end{proposition}

Prescribed-weight vector--functional pairs are a standard tool in these
constructions; compare \cite[Section~3]{ArgyrosLopezAbadTodorcevic2005}
and \cite[Section~1]{ArgyrosArvanitakisTolias2006}.
We prove the estimate needed here directly.

The proof has three steps.  First we find blocks whose norm is spread over
intervals.  Differences of a weak-Cauchy subsequence then permit a Ramsey
argument that limits how often one path can interact with prescribed weight
windows.  Finally we estimate every norming expression by following only
its pieces below the prescribed weight.

\subsection{Blocks and weight windows}

\begin{lemma}\label{pair:partition}
For every integer $N\ge1$, every infinite admissible block sequence has a
finite unit block $u$ such that, for all $1\le d\le N$ and all successive
disjoint ordinal intervals $I_1,\ldots,I_d$,
\[
 \sum_{b=1}^d\|P_{I_b}u\|\le8d^{1/r}.
\]
The block can be chosen in any prescribed tail.
\end{lemma}

\begin{proof}
Work in that tail.  Let $a(k)$ be the supremum of the norms of sums of $k$
successive finite unit blocks.  Contractivity of interval projections gives
$a(k)\le k$.  Norm each of $L_h$ successive unit blocks to within a factor
$2$, crop the normers to their block ranges, and combine them by a Type~I
operation of weight $h$.  In the mixed case use the ordinal and color
ranges.  Thus
\[
 a(L_h)\ge\frac{L_h}{2m_h}.
\]
Choose $t>s$ sufficiently close to $s$ that
$d^{1/t'}\le2d^{1/r}$ for $1\le d\le N$.  The parameter limit in
Lemma~\ref{ch:parameters} implies that $a(k)/k^{1/t}$ is unbounded.
Consequently there is a record index $k$ such that
\[
 a(\ell)\le a(k)(\ell/k)^{1/t}\quad(1\le\ell\le k),
 \qquad a(k)\ge N^{1/s}.
\]
Choose successive unit blocks $v_1,\ldots,v_k$ with
$\|\sum v_i\|\ge a(k)/2$, and normalize their sum to obtain $u$.

For each $I_b$, let $\ell_b$ count the blocks whose ordinal ranges are
contained in $I_b$.  There are at most two further blocks whose ranges
partly meet $I_b$.  Their projected norms are at most $1$ each.  Since
$\sum_b\ell_b\le k$, the record inequality and H\"older's inequality give
\begin{align*}
 \sum_{b=1}^d\|P_{I_b}u\|
 &\le2\sum_{b:\ell_b>0}(\ell_b/k)^{1/t}+\frac{4d}{a(k)}\\
 &\le2d^{1/t'}+\frac{4d}{N^{1/s}}
 \le8d^{1/r}.
\end{align*}
\end{proof}

Fix $J$, put $L=L_J$, and choose $\epsilon>0$ so that
\begin{equation}\label{pair:epsilon}
 \epsilon L^{1/r}\le\theta_J.
\end{equation}
We now construct a weakly null admissible block sequence $(x_i)$ and
integer windows $[J_i,T_i]$ satisfying
\begin{align}
 &1\le\|x_i\|\le2,\qquad J_i>J,
 \qquad J_{i+1}>T_i, \label{pair:window-order}\\
 &J_i>\max\{\rho(\alpha,\beta):
       \alpha\le\beta\text{ belong to }\bigcup_{h<i}\operatorname{supp}x_h\},
       \label{pair:window-coherence}\\
 &\sum_{b=1}^d\|P_{I_b}x_i\|\le A d^{1/r}
       \quad(d\le L_{J_i}),\qquad A=16, \label{pair:window-partition}\\
 &\|x_i\|_1\sum_{j>T_i}m_j^{-1}\le\epsilon,
 \qquad \|x_i\|_1\sum_{j>T_i}m_j^{-1}\longrightarrow0.
       \label{pair:window-tail}
\end{align}
Here $\|x\|_1=\sum_\alpha|x(\alpha)|$ for a finite vector, and the maximum
over an empty set imposes no condition.

For completeness, choose preliminary successive unit blocks $u_i$
recursively.  Choose $J_i>J$, also larger than the preceding $T$ (when
present) and all the $\rho$-values on the supports chosen so far.  Apply
Lemma~\ref{pair:partition} with $N=L_{J_i}$, then choose $T_i\ge J_i$ so
large that
\[
 \|u_i\|_1\sum_{j>T_i}m_j^{-1}\le\min\{\epsilon/2,1/i\}.
\]
By Proposition~\ref{path:weak-cauchy}, pass to a weak-Cauchy subsequence and
take differences of successive pairs.  These differences are weakly null.
Projection onto the first member of a pair gives the lower norm bound $1$;
the upper bound is $2$.  Give each difference the window from the first
member's $J$ to the second member's $T$.  Both old partition bounds apply
for $d\le L_{J_i}$, where $J_i$ is the new lower endpoint, so their sum gives
\eqref{pair:window-partition}.  The remaining assertions follow by
monotonicity of the tails and the original spacing.  We relabel the
resulting sequence and windows as above.

\subsection{A path cannot repeatedly meet its weight windows}

For a finite uncropped path $P=(\phi_\ell)$ with weights $j_\ell$, write
\[
 P[i]=\sum_{\ell:\ J_i\le j_\ell\le T_i}\phi_\ell.
\]
We say that $P$ hits $i$ at threshold $\delta>0$ when
$|P[i](x_i)|>\delta$.  This definition uses the original pieces and their
original coding, even when the path is later cropped inside a norming
expression.

\begin{lemma}\label{pair:four-hits}
For every $\delta>0$ and every infinite set of indices there is an infinite
subset on which no finite uncropped path has four hits at threshold
$\delta$.
\end{lemma}

\begin{proof}
Color a four-element set positively if one finite path hits all four
indices.  By the infinite Ramsey theorem, it suffices to rule out an
infinite set $M$ all of whose four-element subsets are positive.

Fix $v<t<k$ in $M$.  For each $w\in M$ with $w<v$, choose a path hitting
$w,v,t,k$.  Partition these paths according to their actual histories of
weights below $J_v$: this history records the uncropped pieces, their
weights, and their metadata, with their actual ordinal coordinates.
We claim that representatives of different classes have disjoint weight
sets in $[J_k,T_k]$.

Indeed, suppose two such paths share a weight in that window.  It is a
successor weight in both paths, because both hit $t$.  Injectivity of the
coding map identifies their preceding compressed codes.  These codes have
a common final metadata value $p\ge J_t$, so
\eqref{pair:window-coherence} gives
$p>\max\rho(\operatorname{supp}x_v)$.  Choose coordinates $a,b$ of
$\operatorname{supp}x_v$ at which the respective $v$-window pieces are
nonzero, and suppose $a\le b$.  The first closure contains $a$; the second
contains $b$, and hence $a$, by coherence and the bound on $p$.  The two
closures have the same initial segment through $a$, by
Lemma~\ref{ch:closures}.

In the first path every piece of weight below $J_v$ lies strictly before
$a$.  Its compressed array is zero after the position of $a$.  Equality
of the codes therefore identifies that entire uncropped piece in the
second path: the positions through $a$ have the same ordinal labels, and
every later coefficient is zero.  The metadata agree as well.  Thus the
two actual histories below $J_v$ are equal, proving the claim.

Crop each class representative to its pieces in $[J_k,T_k]$.  This is a
legal ordinal crop because pieces and weights are successive.  Their
active weight sets are pairwise disjoint.  Type~II combinations and
rational approximation therefore give
\[
 \big\|(P_b[k](x_k))_b\big\|_{q'}\le\|x_k\|\le2.
\]
Every coordinate has absolute value greater than $\delta$.  In particular,
the number of classes is bounded by an integer $D$ depending only on
$\delta$, for example any integer $D\ge(2/\delta)^{q'}$.

Choose triples $v_n<t_n<k_n$ from $M$ with $v_n\to\infty$, and label the
classes for each triple by elements of $\{1,\ldots,D\}$.  Each fixed
$w\in M$ has a label once $w<v_n$.  Successive subsequences followed by a
diagonal subsequence make its label eventually constant, for every fixed
$w$.  One label occurs for an infinite set $M_0$ of fixed indices.  For
each sufficiently large $n$, let $P_n$ be a representative of that label.
Then, for every fixed $w\in M_0$, $P_n$ hits $w$ for all sufficiently
large $n$: eventually $T_w<J_{v_n}$, and equality of actual histories
below $J_{v_n}$ preserves the whole $w$-window.

Choose $w_0\in M_0$ so large that
\[
 A\sum_{j\ge J_{w_0}}\theta_j<\delta/4.
\]
For fixed $w\in M_0$ with $w>w_0$, consider $n$ large enough that $P_n$
hits both $w_0$ and $w$.  Every piece of weight below $J_{w_0}$ precedes
a coordinate in $\operatorname{supp}x_{w_0}$, and hence vanishes on $x_w$.
For $J_{w_0}\le j<J_w$, the child intervals of a weight-$j$ piece and
\eqref{pair:window-partition} give the bound $A\theta_j$ on its absolute
evaluation at $x_w$.  The total contribution of weights above $T_w$ is at
most the expression in \eqref{pair:window-tail}, which is less than
$\delta/4$ for all sufficiently large $w$.  Consequently
\[
 |P_n(x_w)|\ge\delta/2
\]
for every sufficiently large fixed $w\in M_0$, once $n$ is sufficiently
large depending on $w$.

The finite path sums belong to $B_{E^*}$.  A weak-star convergent subnet of
$(P_n)$ has a limit $P$ satisfying $|P(x_w)|\ge\delta/2$ at every one of
these fixed indices: each inequality is eventually true and defines a
weak-star closed set.  This contradicts weak nullity of $(x_w)$.
\end{proof}

Select $L$ of the blocks successively.  After selecting the $i$th block,
choose $\delta_{i+1}>0$ such that
\begin{equation}\label{pair:threshold}
 T_i\delta_{i+1}\le\epsilon 2^{-i},
\end{equation}
and apply Lemma~\ref{pair:four-hits} to the remaining infinite set.
Relabel the selected blocks $x_1,\ldots,x_L$, carrying their windows with
them.  For each $i<L$, every uncropped path has at most three hits at
threshold $\delta_{i+1}$ among the selected indices strictly after $i$.

\subsection{Estimating a norming expression}

We prove
\begin{equation}\label{pair:sum-bound}
 \left\|\sum_{i=1}^Lx_i\right\|\le500L/m_J.
\end{equation}
In pure mode the relevant norming set is $K_k$; in mixed mode it is $K$,
including its pure terminal functionals.  Lemma~\ref{ch:norm} identifies
the pure norm with the full norm on that color.  The estimate follows by a
recursion on a finite expression in the relevant norming set.  We allow an arbitrary set $I\subset\{1,\ldots,L\}$ to be assigned to each
node; its evaluation vector is $(|f(x_i)|\mathbf1_I(i))_{i=1}^L$.
Allowing these subsets is essential because different descendants will
receive different blocks.

At a nonterminal node, regard Type~I as a single path with coefficient
$1$, and write Type~II as a combination of cropped paths $P_b$ with
coefficients $a_b$, where $\sum_b|a_b|^q\le1$.  Across all these paths the
active weights are distinct.  A group of weight $j$ is a cropped Type~I
piece, represented as $m_j^{-1}$ times a sum of at most $L_j$ nonzero
successive cropped children.  The children have lower construction stage.
The following estimates also apply when some assigned evaluations vanish.
We first bound every evaluation vector in $\ell_r$; we then stop the
same recursion at depth $d_J$ and separate the small $\ell_r$ charges
from the finitely many remaining $\ell_1$ charges.

\paragraph{Weights above a window.}
For each $i\in I$, the total absolute contribution of groups with $j>T_i$
is at most $\epsilon$, by the coefficient bound $1/m_j$ and
\eqref{pair:window-tail}.  Its vector norm in $\ell_r^L$ is at most
$\epsilon L^{1/r}$.

\paragraph{Weights in a window.}
Let $t_{bi}$ be the signed evaluation on $x_i$ of the $b$th cropped path's
pieces with weights in $[J_i,T_i]$.  Further cropping to those pieces
preserves disjoint active weights.  Hence
\begin{equation}\label{pair:column}
 \|(t_{bi})_b\|_{q'}\le2.
\end{equation}
For each path with a nonzero $t_{bi}$ on $I$, let $e$ be the first such
index.  Mark $e$, every subsequent nonzero index at which the path's crop
partly meets $\operatorname{supp}x_i$, and every subsequent index that is
an uncropped hit at threshold $\delta_{e+1}$.  The latter instruction is
empty if $e=L$.  An interval has at most two boundary blocks; a rectangle
has at most four, two for each order.  The thinning gives at most three
large hits after $e$.  Thus each path has at most eight marked indices.
At any other nonzero index the whole block lies in the crop, so its
cropped evaluation agrees with its uncropped evaluation and has absolute
value at most $\delta_{e+1}$.

At most $T_e$ paths can have earliest index $e$, since each then has a
nonzero active piece of a distinct weight in $[J_e,T_e]$.  By
\eqref{pair:threshold}, the total unmarked contribution at each column is
at most $\sum_{e<i}T_e\delta_{e+1}\le\epsilon$.
For the marked contribution $M_i$, H\"older's inequality and
\eqref{pair:column} give
\[
 M_i\le2\left(\sum_{b:\ i\text{ is marked for }b}|a_b|^q\right)^{1/q}.
\]
Summing its $q$th powers and using the eight-mark bound yields
\begin{equation}\label{pair:marked}
 \|M\|_r\le\|M\|_q\le2\,8^{1/q}\le16,
 \qquad \|M\|_1\le16L^{1/q'}.
\end{equation}

\paragraph{Weights below a window.}
For a group of weight $j<J_i$, let $d_{ji}$ be the number of its child
ordinal ranges meeting $\operatorname{supp}x_i$.  If $d_{ji}=0$ there is no
contribution.  Otherwise \eqref{pair:window-partition} applies, because
$d_{ji}\le L_j\le L_{J_i}$, and bounds the absolute contribution by
\[
 \frac{A}{m_j}d_{ji}^{1/r}.
\]
When $d_{ji}=1$, we may instead continue to that unique child, assigning
it $i$ and retaining the coefficient $|a_b|/m_j\le1/m_j$.  Within one
group, the index sets assigned to its different children are disjoint.
When $d_{ji}\ge2$, we stop and charge this bound at the current node.
The boundary incidence estimate is
\begin{equation}\label{pair:incidence}
 \sum_{i:\ d_{ji}\ge2}d_{ji}\le2L_j.
\end{equation}
Indeed, a child range can meet at most one block that also meets an
earlier child range, and at most one that also meets a later child range;
this follows from the succession of both the blocks and the child ranges.
Every counted incidence has one of these two properties.  Thus the
boundary charge for one group has $\ell_r$ norm at most $2A\theta_j$.

\paragraph{A uniform bound on all recursive evaluations.}
If every continued child vector has $\ell_r$ norm at most $K$, disjointness
of the assignments within one group bounds their combined contribution by
$K L_j^{1/r}/m_j=K\theta_j$.  Summing over distinct weights gives $c_0K$.
Assignments from different weights may overlap; this last step uses the
triangle inequality and does not require disjointness between groups.

A terminal coordinate functional has evaluation at most $2$ on at most
one block.  In mixed mode the same is true of a terminal pure functional,
since its color meets at most one mixed block.  Induction on the
construction stage, with arbitrary assigned sets at every stage, now
bounds every node's evaluation vector by $K_0=64$ in $\ell_r$.  In fact
the nonterminal estimate is
\[
 16+2\epsilon L^{1/r}+2A c_0+c_0K_0<64,
\]
by \eqref{pair:epsilon}, $A=16$, $c_0<1/8$, and $\theta_J<1/8$.

\paragraph{Stopping above the prescribed weight.}
We refine this recursion by continuing only the singleton cases with
$j<J$.  Expand these continuations to depth $d=d_J$.  For low groups
with $j\ge J$, stop the singleton cases as well.  Their total incidence
is at most $2L_j+L\le3L_j$, by \eqref{pair:incidence} and $L=L_J\le L_j$.
Their $\ell_r$ charge is consequently at most $3A\theta_j$ per group.
Together with the high and unmarked charges, the resulting small charge
at any node has norm at most
\[
 3A\sum_{j\ge J}\theta_j+2\epsilon L^{1/r}
 \le(12A+2)\theta_J.
\]
Propagate these charges through the continued edges.  The support
condition is preserved at every step: a node assigned $I$ has all its
charge and residual vectors zero outside $I$, and each child receives
only its own assigned subset.  To see why the contraction also applies
to substituted charge vectors, let $w_{jh}\ge0$ be arbitrary vectors
supported on the child assignments of the continued group of weight $j$.
Those supports are disjoint as $h$ varies within a group, although
supports from different groups may overlap.  If every $\|w_{jh}\|_r\le K$,
their recombination, with the path coefficients included, has norm at most
\[
 \sum_{j<J}\frac1{m_j}
       \left(\sum_h\|w_{jh}\|_r^r\right)^{1/r}
 \le K\sum_{j<J}\theta_j\le c_0K.
\]
The first inequality uses disjointness within a group and the triangle
inequality across groups.  Recursive substitution preserves the support
condition, so it costs at most $c_0$ at each level.
Thus all the small charges, together with
the residual evaluation vectors at depth $d$, have total $\ell_r$ norm
at most
\begin{equation}\label{pair:small-total}
 \frac{12A+2}{1-c_0}\theta_J+c_0^dK_0
 \le\left(\frac{194}{7/8}+64\right)\theta_J
 =\frac{2000}{7}\theta_J.
\end{equation}
The depth condition in Lemma~\ref{ch:parameters} was used in the second
inequality.  H\"older's inequality converts this to an $\ell_1$ bound
$\frac{2000}{7}L^{1/s}\theta_J=\frac{2000}{7}L/m_J$.

It remains to count the other charges above depth $d$: terminal charges,
marked charges, and boundary charges of weights $j<J$.  At a node their
$\ell_1$ sum is at most
\[
 2+16L^{1/q'}+A B_J
 \le16(1+B_J+L^{1/q'}).
\]
For the last term we used \eqref{pair:incidence}, $d_{ji}^{1/r}\le d_{ji}$,
and $m_j\ge1$.  A node has at most $\sum_{j<J}L_j$ continued children.
The number of nodes at depths strictly below $d$ is therefore at most
$(1+B_J)^d$, also when $J=1$ and $B_J=0$.  Propagation coefficients are
at most $1$.  The other parameter inequality from
Lemma~\ref{ch:parameters} bounds their total cost by $16L/m_J$.
Combining this with \eqref{pair:small-total} gives
\[
 \sum_{i=1}^L|f(x_i)|\le\frac{2112}{7}\frac L{m_J}
 <500\frac L{m_J}
\]
for every relevant norming functional $f$.  Taking the supremum proves
\eqref{pair:sum-bound}.

\begin{proof}[Proof of Proposition~\ref{pair:exact}]
Perform the preceding construction in the prescribed tail and put
$y=(m_J/L)\sum_{i=1}^Lx_i$.  Equation~\eqref{pair:sum-bound} gives
$\|y\|\le500$.  Choose a normer $g_i$ in the family with
$g_i(x_i)\ge1/2$ and crop it to the ordinal range of $x_i$, also to its
color range in the mixed case.  These nonzero normers are successive, so
\[
 \phi=\frac1{m_J}\sum_{i=1}^Lg_i
\]
is a Type~I functional of weight $J$.  Their cross evaluations vanish,
and hence $\phi(y)=L^{-1}\sum_i g_i(x_i)\ge1/2$.
The crops give the asserted support bounds.
\end{proof}
